\section{sprite Namespace Reference}
\label{namespacesprite}\index{sprite@{sprite}}
The main namespace for sprite.  


\subsection*{Classes}
\begin{CompactItemize}
\item 
struct {\bf sprite::Runtime\-Error}
\begin{CompactList}\small\item\em Thrown to indicate a runtime error in the execution system. \item\end{CompactList}\item 
struct {\bf sprite::Pool\-Item}
\begin{CompactList}\small\item\em An item in the computation pool. \item\end{CompactList}\item 
struct {\bf sprite::Computation\-Pool}
\begin{CompactList}\small\item\em A computation pool. \item\end{CompactList}\item 
struct {\bf sprite::Yield\-Handler}
\begin{CompactList}\small\item\em Base of a class that handles computation results as they are generated. \item\end{CompactList}\item 
struct {\bf sprite::Yield\-Handler\_\-$<$ Output\-Iterator $>$}
\begin{CompactList}\small\item\em A concrete {\bf Yield\-Handler}{\rm (p.\,\pageref{structsprite_1_1YieldHandler})} for a particular output iterator type. \item\end{CompactList}\item 
struct {\bf sprite::Yield\-Handler\_\-$<$ void $>$}
\begin{CompactList}\small\item\em A yield handler that simply discards the output. \item\end{CompactList}\item 
struct {\bf sprite::Fingerprint}
\begin{CompactList}\small\item\em A fingerprint specification. \item\end{CompactList}\item 
struct {\bf sprite::Node}
\begin{CompactList}\small\item\em An expression node. \item\end{CompactList}\item 
struct {\bf sprite::Node\_\-$<$ Payload $>$}
\begin{CompactList}\small\item\em A complete node. The payload contains tag-specific data. \item\end{CompactList}\item 
struct {\bf sprite::Node\-Ptr}
\begin{CompactList}\small\item\em A reference-counting pointer for {\bf Node}{\rm (p.\,\pageref{structsprite_1_1Node})} objects. \item\end{CompactList}\item 
struct {\bf sprite::Program}
\begin{CompactList}\small\item\em A program definition. \item\end{CompactList}\item 
struct {\bf sprite::rewrite$<$ OPER $>$}
\begin{CompactList}\small\item\em Rewrites a node to an OPER node. \item\end{CompactList}\item 
struct {\bf sprite::rewrite$<$ CTOR $>$}
\begin{CompactList}\small\item\em Rewrites a node to a CTOR node. \item\end{CompactList}\item 
struct {\bf sprite::rewrite$<$ FAIL $>$}
\begin{CompactList}\small\item\em Rewrites a node to FAIL. \item\end{CompactList}\item 
struct {\bf sprite::rewrite$<$ CHOICE $>$}
\begin{CompactList}\small\item\em Rewrites a node to CHOICE. \item\end{CompactList}\item 
struct {\bf sprite::rewrite$<$ FWD $>$}
\begin{CompactList}\small\item\em Rewrites a node to FWD. \item\end{CompactList}\item 
struct {\bf sprite::rewrite$<$ INT $>$}
\begin{CompactList}\small\item\em Rewrites a node to INT. \item\end{CompactList}\item 
struct {\bf sprite::rewrite$<$ FLOAT $>$}
\begin{CompactList}\small\item\em Rewrites a node to FLOAT. \item\end{CompactList}\item 
struct {\bf sprite::Loader}
\begin{CompactList}\small\item\em Interface to the module loader. \item\end{CompactList}\item 
struct {\bf sprite::Module}
\begin{CompactList}\small\item\em Interface to a module. \item\end{CompactList}\item 
struct {\bf sprite::User\-Module}
\begin{CompactList}\small\item\em A basic implementation of the {\bf Module}{\rm (p.\,\pageref{structsprite_1_1Module})} interface. \item\end{CompactList}\end{CompactItemize}
\subsection*{Typedefs}
\begin{CompactItemize}
\item 
typedef boost::uint\_\-t$<$ 32 $>$::exact {\bf uint32}\label{namespacesprite_a11}

\item 
typedef {\bf Node\_\-}$<$ {\bf payloads::Unused}$<$$>$ $>$ {\bf Abstract\-Node}\label{namespacesprite_a26}

\begin{CompactList}\small\item\em A complete node where the payload is unknown. \item\end{CompactList}\item 
typedef {\bf Node\_\-}$<$ payloads::In\-Place$<$ 0 $>$ $>$ {\bf In\-Place\-Node0}\label{namespacesprite_a27}

\begin{CompactList}\small\item\em A complete node with exactly zero children. \item\end{CompactList}\item 
typedef {\bf Node\_\-}$<$ payloads::In\-Place$<$ 1 $>$ $>$ {\bf In\-Place\-Node1}\label{namespacesprite_a28}

\begin{CompactList}\small\item\em A complete node with exactly one child. \item\end{CompactList}\item 
typedef {\bf Node\_\-}$<$ payloads::In\-Place$<$ 2 $>$ $>$ {\bf In\-Place\-Node2}\label{namespacesprite_a29}

\begin{CompactList}\small\item\em A complete node with exactly two children. \item\end{CompactList}\item 
typedef {\bf Node\_\-}$<$ {\bf payloads::Child\-List} $>$ {\bf Child\-List\-Node}\label{namespacesprite_a30}

\begin{CompactList}\small\item\em A complete node with more than two children. \item\end{CompactList}\item 
typedef {\bf Node\_\-}$<$ {\bf payloads::Choice} $>$ {\bf Choice\-Node}\label{namespacesprite_a31}

\begin{CompactList}\small\item\em A complete node representing a choice. \item\end{CompactList}\item 
typedef {\bf Node\_\-}$<$ {\bf payloads::Fwd} $>$ {\bf Fwd\-Node}
\begin{CompactList}\small\item\em A complete node representing a forward reference. \item\end{CompactList}\item 
typedef {\bf Node\_\-}$<$ {\bf payloads::Int} $>$ {\bf Int\-Node}\label{namespacesprite_a33}

\begin{CompactList}\small\item\em A complete node representing an integer. \item\end{CompactList}\item 
typedef {\bf Node\_\-}$<$ {\bf payloads::Float} $>$ {\bf Float\-Node}\label{namespacesprite_a34}

\begin{CompactList}\small\item\em A complete node representing a floating-point value. \item\end{CompactList}\item 
typedef tr1::function$<$ void({\bf Node} \&) {\bf h\_\-func\_\-type} )\label{namespacesprite_a43}

\begin{CompactList}\small\item\em The type of an H function. \item\end{CompactList}\end{CompactItemize}
\subsection*{Enumerations}
\begin{CompactItemize}
\item 
enum {\bf Builtin\-Op} \{ \par
{\bf OP\_\-INT\_\-ADD}, 
{\bf OP\_\-INT\_\-SUB}, 
{\bf OP\_\-INT\_\-MUL}, 
{\bf OP\_\-INT\_\-DIV}, 
\par
{\bf OP\_\-FLOAT\_\-ADD}, 
{\bf OP\_\-FLOAT\_\-SUB}, 
{\bf OP\_\-FLOAT\_\-MUL}, 
{\bf OP\_\-FLOAT\_\-DIV}, 
\par
{\bf OP\_\-END}
 \}
\begin{CompactList}\small\item\em The built-in operations. \item\end{CompactList}\item 
enum {\bf Tag\-Value} \{ \par
{\bf FAIL} = 0, 
{\bf OPER}, 
{\bf CHOICE}, 
{\bf FWD}, 
\par
{\bf INT}, 
{\bf FLOAT}, 
{\bf CTOR}
 \}
\begin{CompactList}\small\item\em Identifies the type of a node. CTOR must be last! \item\end{CompactList}\item 
enum {\bf Trace\-Option} \{ {\bf NO\_\-TRACE} = false, 
{\bf TRACE} = true
 \}
\begin{CompactList}\small\item\em Indicates whether tracing is enabled. \item\end{CompactList}\item 
enum {\bf Choice\-Direction} \{ {\bf LEFT} = 0, 
{\bf RIGHT} = 1
 \}
\end{CompactItemize}
\subsection*{Functions}
\begin{CompactItemize}
\item 
template$<$Builtin\-Op Op$>$ void {\bf h\_\-func} ({\bf Node} \&node)\label{namespacesprite_a44}

\begin{CompactList}\small\item\em Implements the H function for a built in operation. \item\end{CompactList}\item 
{\bf Tag\-Value} {\bf make\_\-ctor\_\-tag} (int i)\label{namespacesprite_a45}

\begin{CompactList}\small\item\em Generates a constructor tag. \item\end{CompactList}\item 
bool {\bf is\_\-ctor} ({\bf Tag\-Value} x)\label{namespacesprite_a46}

\begin{CompactList}\small\item\em Indicates whether a tag represents a constructor (built-in or not). \item\end{CompactList}\item 
bool {\bf is\_\-builtin} ({\bf Tag\-Value} x)\label{namespacesprite_a47}

\begin{CompactList}\small\item\em Indicates whether a tag represents a built-in type. \item\end{CompactList}\item 
void {\bf execute} ({\bf Program} const \&pgm, {\bf Node} \&goal, {\bf Yield\-Handler} const \&out, {\bf Trace\-Option} trace)
\begin{CompactList}\small\item\em Execute a goal in the context of the given program. \item\end{CompactList}\item 
void {\bf print\_\-node} ({\bf Node} const \&node)
\begin{CompactList}\small\item\em Print an expression. \item\end{CompactList}\item 
void {\bf execute} ({\bf Program} const \&program, {\bf Node} \&goal, {\bf Trace\-Option} trace)\label{namespacesprite_a51}

\item 
template$<$typename Output\-Iterator$>$ void {\bf execute} ({\bf Program} const \&pgm, {\bf Node} \&goal, Output\-Iterator out, {\bf Trace\-Option} trace=NO\_\-TRACE)\label{namespacesprite_a52}

\begin{CompactList}\small\item\em Alternate form of execute; results are written to the given iterator. \item\end{CompactList}\item 
bool {\bf is\_\-norm} ({\bf Node} const \&node)\label{namespacesprite_a53}

\begin{CompactList}\small\item\em Traverses an expression and returns true if it is found to be in cnf. \item\end{CompactList}\item 
void {\bf pull\_\-tab} ({\bf Node} \&g, {\bf Node\-Ptr} const \&p)
\begin{CompactList}\small\item\em Performs the pull-tab transformation. \item\end{CompactList}\item 
void {\bf head\_\-normalize} ({\bf Node} \&node)
\begin{CompactList}\small\item\em The head-normalizing (H) function from the fair scheme. \item\end{CompactList}\item 
void {\bf fair\_\-normalize} ({\bf Fingerprint} const \&fp, {\bf Node} \&node)\label{namespacesprite_a56}

\begin{CompactList}\small\item\em The normalizing (N) function from the fair scheme. \item\end{CompactList}\item 
{\bf BOOST\_\-STATIC\_\-ASSERT} (sizeof({\bf Abstract\-Node})=={\bf NODE\_\-BYTES})\label{namespacesprite_a57}

\item 
{\bf BOOST\_\-STATIC\_\-ASSERT} (sizeof({\bf In\-Place\-Node0})=={\bf NODE\_\-BYTES})\label{namespacesprite_a58}

\item 
{\bf BOOST\_\-STATIC\_\-ASSERT} (sizeof({\bf In\-Place\-Node1})=={\bf NODE\_\-BYTES})\label{namespacesprite_a59}

\item 
{\bf BOOST\_\-STATIC\_\-ASSERT} (sizeof({\bf In\-Place\-Node2})=={\bf NODE\_\-BYTES})\label{namespacesprite_a60}

\item 
{\bf BOOST\_\-STATIC\_\-ASSERT} (sizeof({\bf Child\-List\-Node})=={\bf NODE\_\-BYTES})\label{namespacesprite_a61}

\item 
{\bf BOOST\_\-STATIC\_\-ASSERT} (sizeof({\bf Choice\-Node})=={\bf NODE\_\-BYTES})\label{namespacesprite_a62}

\item 
{\bf BOOST\_\-STATIC\_\-ASSERT} (sizeof({\bf Fwd\-Node})=={\bf NODE\_\-BYTES})\label{namespacesprite_a63}

\item 
{\bf BOOST\_\-STATIC\_\-ASSERT} (sizeof({\bf Int\-Node})=={\bf NODE\_\-BYTES})\label{namespacesprite_a64}

\item 
{\bf BOOST\_\-STATIC\_\-ASSERT} (sizeof({\bf Float\-Node})=={\bf NODE\_\-BYTES})\label{namespacesprite_a65}

\item 
bool {\bf skip\_\-fwd\_\-ref} ({\bf Node\-Ptr} \&ptr, {\bf Node} \&ref)\label{namespacesprite_a66}

\begin{CompactList}\small\item\em Tells {\bf Node\-Ptr}{\rm (p.\,\pageref{structsprite_1_1NodePtr})} how to skip FWD nodes. \item\end{CompactList}\item 
template$<$typename Visitor$>$ Visitor::result\_\-type {\bf visit} (Visitor const \&visitor, {\bf Node} const \&node)
\begin{CompactList}\small\item\em Unary visitation; calls visitor((Type)(node)). \item\end{CompactList}\item 
template$<$typename Visitor$>$ Visitor::result\_\-type {\bf visit} (Visitor const \&visitor, {\bf Node} \&node)\label{namespacesprite_a68}

\begin{CompactList}\small\item\em Non-const form of visitation. \item\end{CompactList}\item 
template$<$typename Visitor$>$ Visitor::result\_\-type {\bf visit} (Visitor const \&visitor, {\bf Node} const \&lhs, {\bf Node} const \&rhs)
\begin{CompactList}\small\item\em Binary visitation; calls visitor((Type\-Lhs)(lhs), (Type\-Rhs)(rhs)). \item\end{CompactList}\item 
template$<$typename Visitor$>$ Visitor::result\_\-type {\bf visit} (Visitor const \&visitor, {\bf Node} \&lhs, {\bf Node} \&rhs)\label{namespacesprite_a70}

\begin{CompactList}\small\item\em Non-const form of binary visitation. \item\end{CompactList}\end{CompactItemize}
\subsection*{Variables}
\begin{CompactItemize}
\item 
{\bf h\_\-func\_\-type} {\bf builtin\_\-h} [OP\_\-END]
\begin{CompactList}\small\item\em The table of built-in H routines. \item\end{CompactList}\item 
std::string {\bf builtin\_\-oper} [OP\_\-END] = \{ \char`\"{}+\char`\"{}, \char`\"{}-\char`\"{}, \char`\"{}$\ast$\char`\"{}, \char`\"{}/\char`\"{}, \char`\"{}+\char`\"{}, \char`\"{}-\char`\"{}, \char`\"{}$\ast$\char`\"{}, \char`\"{}/\char`\"{} \}\label{namespacesprite_a1}

\begin{CompactList}\small\item\em The tables of built-in operation labels. \item\end{CompactList}\item 
boost::object\_\-pool$<$ {\bf Abstract\-Node} $>$ {\bf node\_\-allocator}\label{namespacesprite_a19}

\begin{CompactList}\small\item\em Global memory pool for allocating nodes. \item\end{CompactList}\item 
{\bf Program} const $\ast$ {\bf g\_\-program} = NULL\label{namespacesprite_a20}

\begin{CompactList}\small\item\em Global pointer to the currently executing program. \item\end{CompactList}\item 
size\_\-t const {\bf NODE\_\-BYTES} = 32
\begin{CompactList}\small\item\em The size of a complete node. \item\end{CompactList}\item 
rewrite$<$ CTOR $>$ const {\bf rewrite\_\-ctor} = rewrite$<$CTOR$>$()\label{namespacesprite_a36}

\begin{CompactList}\small\item\em Rewrites a node to type CTOR. \item\end{CompactList}\item 
rewrite$<$ OPER $>$ const {\bf rewrite\_\-oper} = rewrite$<$OPER$>$()\label{namespacesprite_a37}

\begin{CompactList}\small\item\em Rewrites a node to type OPER. \item\end{CompactList}\item 
rewrite$<$ FAIL $>$ const {\bf rewrite\_\-fail} = rewrite$<$FAIL$>$()\label{namespacesprite_a38}

\begin{CompactList}\small\item\em Rewrites a node to type FAIL. \item\end{CompactList}\item 
rewrite$<$ CHOICE $>$ const {\bf rewrite\_\-choice} = rewrite$<$CHOICE$>$()\label{namespacesprite_a39}

\begin{CompactList}\small\item\em Rewrites a node to type CHOICE. \item\end{CompactList}\item 
rewrite$<$ FWD $>$ const {\bf rewrite\_\-fwd} = rewrite$<$FWD$>$()\label{namespacesprite_a40}

\begin{CompactList}\small\item\em Rewrites a node to type FWD. \item\end{CompactList}\item 
rewrite$<$ INT $>$ const {\bf rewrite\_\-int} = rewrite$<$INT$>$()\label{namespacesprite_a41}

\begin{CompactList}\small\item\em Rewrites a node to type INT. \item\end{CompactList}\item 
rewrite$<$ FLOAT $>$ const {\bf rewrite\_\-float} = rewrite$<$FLOAT$>$()\label{namespacesprite_a42}

\begin{CompactList}\small\item\em Rewrites a node to type FLOAT. \item\end{CompactList}\end{CompactItemize}


\subsection{Detailed Description}
The main namespace for sprite. 

\subsection{Typedef Documentation}
\index{sprite@{sprite}!FwdNode@{FwdNode}}
\index{FwdNode@{FwdNode}!sprite@{sprite}}
\subsubsection{\setlength{\rightskip}{0pt plus 5cm}typedef {\bf Node\_\-}$<${\bf payloads::Fwd}$>$ {\bf sprite::Fwd\-Node}}\label{namespacesprite_a32}


A complete node representing a forward reference. 

This is simply a node that refers to another, which sometimes comes up during rewriting. 

Definition at line 262 of file node.hpp.

\subsection{Enumeration Type Documentation}
\index{sprite@{sprite}!TagValue@{TagValue}}
\index{TagValue@{TagValue}!sprite@{sprite}}
\subsubsection{\setlength{\rightskip}{0pt plus 5cm}enum {\bf Tag\-Value}}\label{namespacesprite_a72}


Identifies the type of a node. CTOR must be last! 

The constructors for any type are assigned the sequential tags CTOR, CTOR+1, CTOR+2,... Said another way, this type is abused in that any value $>$= CTOR can be a valid Tag\-Value. 

Definition at line 66 of file common.hpp.

Referenced by make\_\-ctor\_\-tag(), and sprite::Node::tag().

\subsection{Function Documentation}
\index{sprite@{sprite}!execute@{execute}}
\index{execute@{execute}!sprite@{sprite}}
\subsubsection{\setlength{\rightskip}{0pt plus 5cm}void sprite::execute (Program const \& {\em program}, Node \& {\em goal}, Yield\-Handler const \& {\em handler} = {\tt YieldHandler\_\-$<$\ void\ $>$()}, {\bf Trace\-Option} {\em trace} = {\tt NO\_\-TRACE})}\label{namespacesprite_a49}


Execute a goal in the context of the given program. 

This is the the primary entry point to the execution system. Execution proceeds in-place, meaning that the goal expression will likely be modified.

\begin{Desc}
\item[Parameters:]
\begin{description}
\item[{\em program}]The program definition. This provides the necessary information to carry out execution as a series of rewrite steps. \item[{\em goal}]The goal statement. Modified in place. \item[{\em handler}](optional) Defaults to a handler that discards the generated results. The object that handles output as it is generated. \item[{\em trace}](optional) Defaults to true. Enables debug tracing. \end{description}
\end{Desc}


Definition at line 102 of file exec.cpp.

References sprite::Fingerprint::at(), fair\_\-normalize(), g\_\-program, sprite::Fingerprint::has(), sprite::Node::id(), is\_\-norm(), sprite::Fingerprint::resize(), sprite::Fingerprint::set(), sprite::Fingerprint::size(), sprite::Node::tag(), and sprite::Yield\-Handler::yield().

Here is the call graph for this function:\begin{figure}[H]
\begin{center}
\leavevmode
\includegraphics[width=282pt]{namespacesprite_a49_cgraph}
\end{center}
\end{figure}
\index{sprite@{sprite}!head_normalize@{head\_\-normalize}}
\index{head_normalize@{head\_\-normalize}!sprite@{sprite}}
\subsubsection{\setlength{\rightskip}{0pt plus 5cm}void head\_\-normalize (Node \& {\em node})\hspace{0.3cm}{\tt  [inline]}}\label{namespacesprite_a55}


The head-normalizing (H) function from the fair scheme. 

This is a generic H function that dispatches to the correct implementation for the given operation. The H.6 rule, which ignors constructor-rooted expressions, is implemented here so that user-compiled H rules can ignore it. 

Definition at line 139 of file exec.hpp.

References g\_\-program, sprite::Node::id(), sprite::Program::oper, and sprite::Node::tag().

Referenced by fair\_\-normalize(), and h\_\-func().

Here is the call graph for this function:\begin{figure}[H]
\begin{center}
\leavevmode
\includegraphics[width=140pt]{namespacesprite_a55_cgraph}
\end{center}
\end{figure}
\index{sprite@{sprite}!print_node@{print\_\-node}}
\index{print_node@{print\_\-node}!sprite@{sprite}}
\subsubsection{\setlength{\rightskip}{0pt plus 5cm}void sprite::print\_\-node (Node const \&)}\label{namespacesprite_a50}


Print an expression. 

\begin{Desc}
\item[Note:]Alias for operator$<$$<$, callable from the debugger. \end{Desc}


Definition at line 190 of file exec.cpp.\index{sprite@{sprite}!pull_tab@{pull\_\-tab}}
\index{pull_tab@{pull\_\-tab}!sprite@{sprite}}
\subsubsection{\setlength{\rightskip}{0pt plus 5cm}void pull\_\-tab (Node \& {\em g}, Node\-Ptr const \& {\em p})\hspace{0.3cm}{\tt  [inline]}}\label{namespacesprite_a54}


Performs the pull-tab transformation. 

\begin{Desc}
\item[Parameters:]
\begin{description}
\item[{\em g}]The source of the transformation. \item[{\em p}]The target of the transformation. Must be a direct descendant of g. \end{description}
\end{Desc}


Definition at line 120 of file exec.hpp.

References sprite::Node::clone(), sprite::Node::position(), and rewrite\_\-choice.

Referenced by fair\_\-normalize(), and h\_\-func().

Here is the call graph for this function:\begin{figure}[H]
\begin{center}
\leavevmode
\includegraphics[width=248pt]{namespacesprite_a54_cgraph}
\end{center}
\end{figure}
\index{sprite@{sprite}!visit@{visit}}
\index{visit@{visit}!sprite@{sprite}}
\subsubsection{\setlength{\rightskip}{0pt plus 5cm}template$<$typename Visitor$>$ Visitor::result\_\-type visit (Visitor const \& {\em visitor}, Node const \& {\em lhs}, Node const \& {\em rhs})\hspace{0.3cm}{\tt  [inline]}}\label{namespacesprite_a69}


Binary visitation; calls visitor((Type\-Lhs)(lhs), (Type\-Rhs)(rhs)). 

Determines the type of each node, then calls the visitor with both nodes cast to those types.

Rather than support all four possible overloads related to the const qualifier on {\bf Node}{\rm (p.\,\pageref{structsprite_1_1Node})} arguments, just support the cases where both are qualified or neither one is. 

Definition at line 142 of file visit.hpp.\index{sprite@{sprite}!visit@{visit}}
\index{visit@{visit}!sprite@{sprite}}
\subsubsection{\setlength{\rightskip}{0pt plus 5cm}template$<$typename Visitor$>$ Visitor::result\_\-type visit (Visitor const \& {\em visitor}, Node const \& {\em node})\hspace{0.3cm}{\tt  [inline]}}\label{namespacesprite_a67}


Unary visitation; calls visitor((Type)(node)). 

Determines the type of the node, based on its tag and arity, then calls the visitor with the node cast to that type. 

Definition at line 18 of file visit.hpp.

References sprite::Node::arity(), and sprite::Node::tag().

Referenced by sprite::Node::clone(), sprite::Node::iter(), sprite::Node::operator[$\,$](), and sprite::Node::position().

Here is the call graph for this function:\begin{figure}[H]
\begin{center}
\leavevmode
\includegraphics[width=116pt]{namespacesprite_a67_cgraph}
\end{center}
\end{figure}


\subsection{Variable Documentation}
\index{sprite@{sprite}!builtin_h@{builtin\_\-h}}
\index{builtin_h@{builtin\_\-h}!sprite@{sprite}}
\subsubsection{\setlength{\rightskip}{0pt plus 5cm}{\bf h\_\-func\_\-type} {\bf sprite::builtin\_\-h}}\label{namespacesprite_a0}


{\bf Initial value:}

\footnotesize\begin{verbatim}
  {
      &h_func<OP_INT_ADD>
    , &h_func<OP_INT_SUB>
    , &h_func<OP_INT_MUL>
    , &h_func<OP_INT_DIV>
    , &h_func<OP_FLOAT_ADD>
    , &h_func<OP_FLOAT_SUB>
    , &h_func<OP_FLOAT_MUL>
    , &h_func<OP_FLOAT_DIV>
  }
\end{verbatim}\normalsize 
The table of built-in H routines. 



Definition at line 189 of file builtins.cpp.\index{sprite@{sprite}!NODE_BYTES@{NODE\_\-BYTES}}
\index{NODE_BYTES@{NODE\_\-BYTES}!sprite@{sprite}}
\subsubsection{\setlength{\rightskip}{0pt plus 5cm}size\_\-t const {\bf sprite::NODE\_\-BYTES} = 32}\label{namespacesprite_a25}


The size of a complete node. 

Constrining every complete node type to have the same size simplifies memory pooling and rewriting. 

Definition at line 228 of file node.hpp.